\documentclass[10pt,a4paper]{amsart}
\usepackage[margin=19mm]{geometry}
\usepackage{amsmath,amssymb,mathtools}
\usepackage[scr=rsfs,cal=euler]{mathalfa}
\usepackage{xfrac}
\usepackage[unicode,hidelinks,bookmarks=false]{hyperref}
\newcommand{\ie}{i.e.\ }
\newcommand{\cf}{cf.\ }
\newcommand{\resp}{respectively\ }
\newcommand{\Subsection}{\subsection}
\providecommand{\nobreakdash}{\nobreak\hbox{-}}
\setlength{\emergencystretch}{2em}
\multlinegap0pt
\usepackage{amssymb}
\usepackage{tikz}
\usepackage{xcolor}
\newtheorem{lem}{Lemma}
\newcommand{\GL}{\mathrm{GL}}
\newcommand{\RR}{{\mathbb R}}
\newcommand{\tr}{\mathrm{tr}\,}
\title{Non-linear geometry of multiple zeta values}
\author{Francis Brown}
\date{Source: arXiv:2604.22735v1 (CC BY 4.0)}
\begin{document}
\maketitle

\noindent\emph{Source and licence.} Non-linear geometry of multiple zeta values. Version arXiv:2604.22735v1 (CC BY 4.0), distributed by arXiv under the Creative Commons Attribution 4.0 International licence, http://creativecommons.org/licenses/by/4.0/, as recorded in the arXiv licence metadata for this exact version and shown on the abstract page https://arxiv.org/abs/2604.22735v1. Full licence notice: \texttt{licenses/arxiv-2604.22735-CC-BY-4.0.txt}.\par\smallskip

\noindent\textit{Editorial scope.} Counted unit: the article's lem lem (source0.tex lines 1778--1780). The article's complete original proof of it is delivered with it (source0.tex lines 1781--1786, one \texttt{proof} environment of 6 lines). No mathematics is written, repaired or re-derived: the statement and the proof are the article's own lines, in the article's own notation, and no retained line is substituted or re-flowed.  No further source-local context is needed: every cross-reference of every retained line resolves inside the two delivered spans.  Nothing was omitted from any retained span, so the package carries no omission record.  No retained line contains a citation, so the wrapper carries no bibliography and no bibliography entry.  No retained line contains a figure environment, an image-inclusion command or drawing code, so no image is delivered and the package declares no asset dependency.  Editorial formatting only, in addition: an AMS \texttt{amsart} preamble that restates the article's own declarations and macros used by the retained lines, namely the environments \texttt{lem} on separate counters, together with the standard packages those lines need.  The retained environments print in this document as Lemma 1 in order of appearance, whereas the article prints the same items with the numbering of its own full document.  The complete native source of the article as distributed in the arXiv e-print for this exact version is bundled unchanged as \texttt{sources/199123/source0.tex}.
\begin{lem} Let $X$ be an invertible  $m\times m$ matrix as above.  For any invertible matrices $A, B \in \GL_m(\RR)$ with constant entries (i.e. $dA= dB=0$), one has
\[ \omega^n_{X}= \omega^n_{AX} = \omega^n_{XB}\ .\] 
\end{lem} 

\begin{proof} The form $X^{-1} dX$ is already invariant by left-multiplication by $A$. Indeed, 
\[  (AX)^{-1} d(AX) =  X A^{-1} A \, dX = X^{-1} dX\ .\]
It follows that $\omega^n_X$ is also left-invariant. Now observe that  $dX X^{-1}$ is invariant by right-multiplication by $B$ and notice that
\[  \omega^n_X = \tr (X^{-1} dX. \ldots X^{-1} dX) = \tr (  dX X^{-1}. \ldots dX X^{-1})\]
by cyclicity of the trace.  
\end{proof} 

\end{document}
